\chapter{希尔伯特谱理论}
\setcounter{exercisepar}{0}

对每个向量空间 E，我们用 \(1_{E}\) 表示 E 的恒等映射。若 E、F 与 G 是向量空间，且 v: \(E \rightarrow F\) 与 u: \(F \rightarrow G\) 是线性映射，则从 E 到 G 的线性映射 \(u \circ v\) 有时记作 uv。

设 E 是一个预希尔伯特空间。我们用 \(\langle x|y\rangle_{E}\) ，或简作 \(\langle x|y\rangle\) ，表示 E 上的内积。E 的子集 A 的正交记作 A°。它是 E 的一个闭子空间；它化为 0 当且仅当 A 在 E 中是全的（EVT, V, p. 16, 推论 1）。若 F ⊂ E 是 E 的向量子空间，则 F°° = \(\overline{F}\) （参见 EVT, V, p. 13）。E 的自同态的谱总是相对于巴拿赫代数 \(\mathcal{L}(E)\) 而言的。

设 F 是希尔伯特空间。我们用 \(\mathcal{L}_1(\mathrm{E})\) 表示 E 的有限迹自同态所成的空间（EVT, V, p. 50, 定义 8），用 \(\mathcal{L}_2(\mathrm{E};\mathrm{F})\) 表示从 E 到 F 的希尔伯特–施密特映射所成的空间（EVT, V, p. 51, 定义 9）。映射 \(u\mapsto \| u\| _2 = (\mathrm{Tr}(u^{*}u))^{1 / 2}\) 是 \(\mathcal{L}_2(\mathrm{E};\mathrm{F})\) 上的一个希尔伯特范数（EVT, V, p. 52, 定理 1, (i)）。

有限秩的连续线性映射所成的空间 \(\mathcal{L}^{\mathrm{f}}(\mathrm{E};\mathrm{F})\) 在 \(\mathcal{L}^{\mathrm{c}}(\mathrm{E};\mathrm{F})\) 中稠密（III, p. 15, 命题 10 与 III, p. 16, 命题 11 的推论）。特别地，空间 \(\mathcal{L}_2(\mathrm{E};\mathrm{F})\) 在 \(\mathcal{L}^{\mathrm{c}}(\mathrm{E};\mathrm{F})\) 中稠密。

E 与 F 的希尔伯特张量积（EVT, V, p. 28, 定义 1）记作 E \(\widehat{\otimes}_2\) F。

若 X 是局部紧拓扑空间，\(\mu\) 是 X 上的一个测度，我们记 \(\mathcal{L}^{p}(\mathrm{X},\mu)=\mathcal{L}_{\mathbf{C}}^{p}(\mathrm{X},\mu)\) 与 \(\mathrm{L}^{p}(\mathrm{X},\mu)=\mathrm{L}_{\mathbf{C}}^{p}(\mathrm{X},\mu)\) （对 \(p\in[1,+\infty]\) ）。

